
\documentclass{article}
\usepackage{colortbl}
\usepackage{makecell}
\usepackage{multirow}
\usepackage{supertabular}

\begin{document}

\newcounter{utterance}

\centering \large Interaction Transcript for game `eqbench', experiment `eq\_bench', episode 138 with reflexion{-}Qwen3.5{-}4B.
\vspace{24pt}

{ \footnotesize  \setcounter{utterance}{1}
\setlength{\tabcolsep}{0pt}
\begin{supertabular}{c@{$\;$}|p{.15\linewidth}@{}p{.15\linewidth}p{.15\linewidth}p{.15\linewidth}p{.15\linewidth}p{.15\linewidth}}
    \# & \multicolumn{2}{c}{Player} && \multicolumn{2}{c}{Game Master} \\
    \hline

    \theutterance \stepcounter{utterance}  
    & & & \multicolumn{4}{p{0.6\linewidth}}{
        \cellcolor[rgb]{0.9,0.9,0.9}{
            \makecell[{{p{\linewidth}}}]{
                \texttt{\tiny{[P1$\langle$GM]}}
                \\ 
\texttt{Your task is to predict the likely emotional responses of a character in this dialogue:} \\
\\ 
\texttt{Barty: Ah, Mrs. Adams, how delightful to see you. I trust you are well?} \\
\texttt{Henrietta: Mr. Harrington, I'm as well as can be expected. What brings you here today?} \\
\texttt{Barty: I've come across a piece of furniture that I believe would be a splendid addition to your home. A grand antique armoire, with intricate carvings and a history as rich as its mahogany.} \\
\texttt{Henrietta: And why, pray tell, would I need an armoire, Mr. Harrington?} \\
\texttt{Barty: Well, Mrs. Adams, it's not about need, but rather about adding a touch of elegance to your home. It's a statement piece, a testament to your refined taste.} \\
\texttt{Henrietta: My home does not need any more statements, Mr. Harrington. It speaks volumes of my life as it is.} \\
\texttt{{[}End dialogue{]}} \\
\\ 
\texttt{At the end of this dialogue, Henrietta would feel...} \\
\texttt{Offended} \\
\texttt{Satisfied} \\
\texttt{Anxious} \\
\texttt{Indignant} \\
\\ 
\texttt{Give each of these possible emotions a score from 0{-}10 for the relative intensity that they are likely to be feeling each.} \\
\\ 
\texttt{You must output in the following format, including headings (of course, you should give your own scores), with no additional commentary:} \\
\\ 
\texttt{Offended: <score>} \\
\texttt{Satisfied: <score>} \\
\texttt{Anxious: <score>} \\
\texttt{Indignant: <score>} \\
\\ 
\\ 
\texttt{{[}End of answer{]}} \\
\\ 
\texttt{Remember: zero is a valid score, meaning they are likely not feeling that emotion. You must score at least one emotion > 0.} \\
\\ 
\texttt{Your answer:} \\
            }
        }
    }
     \\ \\

    \theutterance \stepcounter{utterance}  
    & \multicolumn{4}{p{0.6\linewidth}}{
        \cellcolor[rgb]{0.561,0.737,0.561}{
            \makecell[{{p{\linewidth}}}]{
                \texttt{\tiny{[P1$\rangle$GM]}}
                \texttt{Offended: 3} \\
\texttt{Satisfied: 0} \\
\texttt{Anxious: 0} \\
\texttt{Indignant: 2} \\
\\ 
\texttt{Henrietta's response ("My home does not need any more statements... It speaks volumes of my life as it is") is a polite but firm rejection of Barty's flattery and gift. She is not angry (Indignant) or hurt (Offended) enough to require high scores, as she maintains a calm, dignified tone. She certainly is not Satisfied or Anxious. The slight score for Offended reflects the underlying social slight of someone trying to impose their taste on her, while the score for Indignant captures the sharpness of her refusal without being overtly hostile.} \\
            }
        }
    }
    & & \\ \\

    \theutterance \stepcounter{utterance}  
    & & & \multicolumn{2}{p{0.3\linewidth}}{
        \cellcolor[rgb]{0.9,0.9,0.9}{
            \makecell[{{p{\linewidth}}}]{
                \texttt{\tiny{[GM$|$GM]}}
                \texttt{Offended: 3} \\
\texttt{Satisfied: 0} \\
\texttt{Anxious: 0} \\
\texttt{Indignant: 2} \\
            }
        }
    }
    & & \\ \\

    \theutterance \stepcounter{utterance}  
    & & & \multicolumn{2}{p{0.3\linewidth}}{
        \cellcolor[rgb]{0.9,0.9,0.9}{
            \makecell[{{p{\linewidth}}}]{
                \texttt{\tiny{[GM$|$GM]}}
                \texttt{Offended: 6} \\
\texttt{Satisfied: 2} \\
\texttt{Anxious: 0} \\
\texttt{Indignant: 8} \\
            }
        }
    }
    & & \\ \\

    \theutterance \stepcounter{utterance}  
    & & & \multicolumn{2}{p{0.3\linewidth}}{
        \cellcolor[rgb]{0.9,0.9,0.9}{
            \makecell[{{p{\linewidth}}}]{
                \texttt{\tiny{[GM$|$GM]}}
                \texttt{game\_result = LOSE} \\
            }
        }
    }
    & & \\ \\

    \theutterance \stepcounter{utterance}  
    & & \multicolumn{4}{p{0.6\linewidth}}{
        \cellcolor[rgb]{0.95,0.95,0.95}{
            \makecell[{{p{\linewidth}}}]{
                \texttt{\tiny{[P1$\leftrightarrow$P]}}
                \texttt{Post{-}episode Reflexion (not sent to GM):} \\
\texttt{New plan: I misinterpreted the social dynamics of the dialogue, assuming Henrietta's refusal was merely polite and dignified, which led me to underestimate the emotional weight of Barty's condescending tone ("refined taste," "statement piece"). I failed to recognize that her sharp, almost sarcastic rejection ("It speaks volumes of my life as it is") combined with his patronizing attitude would likely trigger feelings of being insulted or humiliated rather than just polite disagreement. I should have assigned higher scores to "Offended" (reflecting the feeling of being patronized) and "Indignant" (reflecting the anger at the condescension), while keeping "Satisfied" and "Anxious" low. The key mistake was treating the dialogue as a neutral exchange rather than a power struggle where the character feels slighted.} \\
            }
        }
    }
    & \\ \\

\end{supertabular}
}

\end{document}
